% lit_membership.tex -- Literature 2.2: changes of team membership, with the numeric table T3 (drafted 12 Oct 2026
% from the thesis plan of 11 Oct). The facts on Quan et al. and Zeng et al. were verified from the full texts on
% 11 Oct (hand-back docs/thesis_work_2026-10-11.md); the cells in brackets wait for a full-text read (plan C2-A:
% Bacheti and Delbene were read from abstracts only) or for the October flights. Entries the plan lists for this
% section that are not yet in the bibliography (Oishi 2023, Gao 2024, Delbene 2024/2025, Bacheti 2026, Xu 2025) wait for the approval batch;
% Fink 2011, Tognon 2018 and Farina 2020 were added 12 Oct.
% Papers found with Undermind and approved by Wesley, 10 Oct (DOIs, abstracts and sources in
% docs/thesis_drafts/refs_2026-10-10.md; alternates listed there).

\subsection{Changes of team membership}
\label{sec:lit_team}

In this thesis a reconfiguration is a change of team membership: a vehicle leaving or joining while the payload
is carried. The word is also used for a fixed team changing its shape, as when a team contracts to pass a doorway
\cite{jackson2020scalable, tensionbased}; that is a formation change, and the team's force allocation keeps its
dimension. A membership change does not. Most of the work of Section~\ref{sec:lit_models} assumes a fixed team
\cite{efficient, wahbaKinodynamicMotionPlanning2024, quan2025swarm}; the few works that change it are set out in
Table~\ref{tab:membership}, with the detection of the change kept apart from the recovery from it, because the
two are often reported as one time.

\paragraph{A vehicle leaving, on hardware.} Quan et al.\ \cite{quan2025swarm} remove a vehicle from a carried
load in flight: four robots carry a 4\,kg load, one releases its cable on command, and the other three
reorganise to a near-equilateral layout and rebalance their tensions within 13\,s, decentralised and without
replanning. They report the cable tensions and one robot's velocity through the event, not the load's attitude.
Zeng et al.\ \cite{zeng2025marl} switch one of three vehicles off in flight and leave it hanging from a 1.4\,kg
load; the other two keep five degrees of freedom of the load under a learned decentralised policy that observes
only the load's pose, and the load's position and attitude are plotted through the loss with no transient
number. Bacheti et al.\ \cite{tensionbased} handle a removal by repositioning the remaining vehicles before the
release so that the transient is small, on hardware with a point-mass payload; [whether the agent count changes
in their experiments is to be verified from the full text]. So on hardware a vehicle has been released from a
carried load and a vehicle has been lost from one, and in neither case is the load's attitude through the event
given as a number.

\paragraph{A link that fails, in simulation.} A cable that breaks rather than releases is treated as a fault: the
remaining tensions are estimated and the load reassigned \cite{liang2021fault}, a supervisor reconfigures a
passivity-based formation \cite{saiella2021fault}, or each vehicle's measured tension is routed into its thrust so
that five vehicles recover from one or two severed cables within a pendulum period \cite{hajieghrary2026passive}; see
also \cite{farrasFaultTolerantControlMultiQuadcopter2024}. None flies, where a time is given it is the recovery and not
the detection, and none adds a vehicle.

\paragraph{A vehicle joining.} No work attaches a vehicle to a payload that a team is already carrying. Quan et al.\
describe a joining robot only in prose \cite{quan2025swarm}, and the event-triggered architecture of Tasooji et al.\
\cite{tasoojiCooperativeControlMultiQuadrotors} is argued to scale but not analysed for a varying team. The closest
hardware precedent is a payload handed from one multirotor to another, the grasp confirmed by a magnetic sensor before
the giver releases \cite{shankar2020inflight}; the team does not grow and the payload is not shared.

\paragraph{Prepared changes and statics.} Plug-and-play MPC lets subsystems join or leave a running network without
losing stability, provided each change is prepared before it happens \cite{riverso2013plug, zeilinger2013plug,
aboudonia2022reconfigurable, saccani2023model}, all on linear coupled subsystems or ground robots; a suspended payload
couples every vehicle through the load, and a magnet that lets go happens at an instant the planner does not choose, so
this thesis switches to a pre-built problem of the new size at the event and shows the first solve in flight rather
than certifying it. Static equilibrium decides which layouts can hold a payload at all: three robots are the fewest,
their attachment points must span its centre of mass \cite{michael2011cooperative, jiang2013inverse}, and the shares
are exact when the centre lies in the attachment plane \cite{tognon2018internal}; see also
\cite{gabellieri2023equilibria, liu2025wrench, fink2011planning}. All analyse a fixed team; Section~\ref{sec:reconfig}
turns the condition into a margin that says which carrier a layout may lose.

\begin{table}[htbp]
\centering
\footnotesize
\setlength{\tabcolsep}{3pt}
\caption[Changes of team membership under a carried payload]{Changes of team membership under a carried
payload, as each work reports them. Announced: the controller is told before it happens. Leaver: whether the vehicle
leaves the load or stays attached. Detection: from the event to the controller acting on it; recovery: to a settled
state. ``--'': not reported; brackets: to fill from the full texts.}
\label{tab:membership}
\begin{tabularx}{\textwidth}{@{}
  >{\raggedright\arraybackslash\hsize=0.86\hsize}X
  >{\raggedright\arraybackslash\hsize=1.04\hsize}X
  >{\raggedright\arraybackslash\hsize=0.78\hsize}X
  >{\raggedright\arraybackslash\hsize=0.90\hsize}X
  >{\raggedright\arraybackslash\hsize=1.04\hsize}X
  >{\raggedright\arraybackslash\hsize=1.14\hsize}X
  >{\raggedright\arraybackslash\hsize=1.24\hsize}X@{}}
\toprule
Work & Event; mechanism & Announced & Leaver & Plant; team; payload & Attitude or sag through the event &
Detection; recovery \\
\midrule
Quan\newline \cite{quan2025swarm}, 2025 & release on command; cable released at the load & yes & removed, flies off with
its cable & hardware, outdoor; 4 to 3; 4\,kg & -- (tensions and one robot's velocity plotted) & --; tensions
balanced after 13\,s \\
\addlinespace
Zeng\newline \cite{zeng2025marl}, 2025 & motors switched off & no & left hanging from the load & hardware; 3 to 2;
1.4\,kg & plotted, no number & none explicit (learned policy); -- \\
\addlinespace
Bacheti\newline \cite{tensionbased}, 2023 & planned release; the team repositions first & yes & [removed; count to
verify] & hardware; [3 to 2 to 3]; point mass & -- & --; -- \\
\addlinespace
Liang\newline \cite{liang2021fault}, 2021 & cable breaks & no & removed & simulation; [team]; [mass] & [from the paper] &
disturbance observers; [from the paper] \\
\addlinespace
Saiella\newline \cite{saiella2021fault}, 2021 & vehicle fault; the supervisor reconfigures & no & [from the paper] &
simulation; [team] & -- & supervisor; -- \\
\addlinespace
Hajieghrary\newline \cite{hajieghrary2026passive}, 2026 (preprint) & one or two cables severed & no & removed &
simulation; 5 to 4 and 5 to 3 & [from the preprint] & none (passive feed-forward); within a pendulum period
[0.32\,s] \\
\addlinespace
Farras\newline \cite{farrasFaultTolerantControlMultiQuadcopter2024}, 2024 & [fault, from the paper] & no & [from the
paper] & simulation; [team] & [from the paper] & [from the paper] \\
\addlinespace
Shankar\newline \cite{shankar2020inflight}, 2020 & a bar handed from one multirotor to another & yes & the giver releases
after the grasp is confirmed & hardware; 1 to 1 & -- & magnetic sensor at the grasp; gains rescheduled \\
\addlinespace
ModQuad\newline \cite{saldana2018modquad}, TRADY \cite{sugihara2023trady}, LEGION \cite{sugihara2026legion}, Miyazaki \cite{miyazaki2018airborne} & vehicles docking to one another; no payload & yes & n/a & hardware; 1 to 1 per dock &
TRADY: altitude jump 0.39\,m, 0.065\,m with the thrust blend & motion capture, contact force or a latch; -- \\
\midrule
This work & magnet off, commanded or not (leave); contact and co-motion (join) & both & removed; lands by itself
& hardware and twin; 4 to 3, 3 to 4 (twin); 0.86\,kg ring on rods & tilt rise 5.3--16.9$^\circ$ over five rig
releases & 0.33--0.42\,s (rig); join: [join tree] \\
\bottomrule
\end{tabularx}
\end{table}
